% titel: Kanal
% skript: Kap. Hydrodynamik

% Eine leicht geneigte, offene Ölrinnen (sehr breit) wird von einer
% gravitationsgetriebenen Schichtenströmung durchflossen. Gefragt sind die
% Bedeutung der Reynoldszahl, das Oberflächen-Geschwindigkeitsprofil und die
% Reynoldszahl der Strömung (laminar oder turbulent?).

Eine sehr breite, leicht geneigte Rinne der Tiefe $d$ wird von Öl durchströmt
(seitliche Reibung vernachlässigbar). Gesucht sind die Charakteristik der
Reynoldszahl, das Geschwindigkeitsprofil an der Oberfläche und die
Reynoldszahl der Strömung.

\begin{werkzeugkasten}[Werkzeuge]
\textbf{Reynoldszahl:} dimensionsloses Verhältnis von Trägheits- zu
Viskosenkräften,
\begin{equation}
  \mathrm{Re}=\frac{\rho\, \bar v\, L}{\eta}\;,
\end{equation}
$\bar v$ = mittlere Geschwindigkeit, $L$ = charakteristische Länge (hier die
Tiefe $d$). Sie entscheidet, ob die Strömung laminar oder turbulent ist
(laminar für kleine $\mathrm{Re}$, turbulent ab $\sim$Kritischer Zahl).

\textbf{Geschwindigkeitsprofil einer Schichtenströmung:} Die Newtonsche
Viskositätsbeziehung $F_R=\eta A\,\mathrm dv/\mathrm dy$ (Schubspannung
$\tau=\eta\,\mathrm dv/\mathrm dy$) wird mit den antreibenden Kräften
gleichgesetzt. Bei geneigter, gravitationsgetriebener Strömung treibt die
Hangabtriebskomponente $\rho g A (d-y)\sin\alpha$ die darüberliegende Schicht
an, was nach Integration ein parabelförmiges Profil liefert,
\begin{equation}
  v(y)=\frac{\rho g\sin\alpha}{\eta}\Big(d\,y-\frac{y^2}{2}\Big)\;,
  \qquad \text{mit } v(0)=0\ \Rightarrow\
  v(d)=\frac{\rho g\sin\alpha\, d^2}{2\eta}\;.
\end{equation}

\textbf{Mittlere Geschwindigkeit (Parabelprofil):}
\begin{equation}
  \bar v=\frac{1}{d}\int_0^d v(y)\,\mathrm dy
  =\frac{2}{3}\,v(d)\;.
\end{equation}

\textbf{Einheiten:} $\eta$ in \unit{\pascal\second}, $\rho$ in
\unit{\kilogram\per\cubic\metre}, $g\approx\qty{9.81}{m/s^2}$; $\sin\alpha$
dimensionslos.
\end{werkzeugkasten}

\begin{werkzeugkasten}[Vorgehen in allgemeiner Form]
\begin{enumerate}
  \item (a) Reynoldszahl definieren und ihre Bedeutung benennen.
  \item (b) Schubspannungsgleichgewicht für eine lamellare Schicht
       aufstellen, nach $\mathrm dv/\mathrm dy$ auflösen, mit $v(0)=0$
       integrieren --- Ergebnis ist das Profil, an der Oberfläche $y=d$.
  \item (c) Mittlere (über die Tiefe gemittelte) Geschwindigkeit bilden, in
       $\mathrm{Re}$ einsetzen, mit der kritischen Zahl vergleichen.
\end{enumerate}
\end{werkzeugkasten}

\begin{achtung}[Typische Fallen]
\begin{itemize}
  \item $\sin\alpha$ aus der Höhendifferenz $\Delta h$ über die Länge $L$:
       $\sin\alpha\approx \Delta h/L$ (kleine Neigung).
  \item Mittlere Geschwindigkeit $\neq$ Profil an der Oberfläche: für ein
       Parabelprofil ist $\bar v=\frac23 v(d)$.
  \item Viskosität umrechnen ($\mathrm{mPa\,s}\to\mathrm{Pa\,s}$); Dichte in
       $\mathrm{kg/m^3}$ (nicht $\mathrm{kg/l}$).
  \item $d$ als charakteristische Länge, nicht die Rinnenbreite.
\end{itemize}
\end{achtung}
